% !TEX root = ../Main.tex
\chapter{二次方程的韦达定理}

\section{韦达定理}

对一元二次方程
\[
ax^2+bx+c=0\quad(a\ne 0),
\]
\textbf{使用前提}是方程有实根，即 $\Delta=b^2-4ac\ge 0$。

\state{韦达定理}{
    设两实根为 $x_1,x_2$，则
    \[
    \bc
    x_1+x_2=-\dfrac{b}{a},\\[4pt]
    x_1x_2=\dfrac{c}{a}.
    \ec
    \]
}

由"和"与"积"可以立刻拼出"根的差"的绝对值：
\[
\abs{x_1-x_2}=\sqrt{(x_1+x_2)^2-4x_1x_2}
=\sqrt{\frac{b^2-4ac}{a^2}}
=\frac{\sqrt{\Delta}}{\abs{a}}.
\]

\rmkb{
    韦达定理体现的思想是：$x_1$ 与 $x_2$ 的\textbf{和、积、差}这些对称关系，可以直接从系数读出来，\textbf{不必先把 $x_1,x_2$ 各自解出来}。所以做题时——能不具体求根，就不求。
}

\section{弦长公式}

直线与曲线交于 $A(x_1,y_1)$、$B(x_2,y_2)$，要算弦长 $\abs{AB}$。设直线斜率 $k=\tan\theta$（$\theta$ 为倾斜角）。

\begin{center}
\begin{tikzpicture}[>=Stealth,scale=1.2]
  \coordinate (A) at (0,0);
  \coordinate (B) at (2.4,1.4);
  \draw[thick,blue!60!black] (-0.3,-0.175) -- (2.8,1.633);
  \draw[dashed] (A) -- (B |- A) node[midway,below]{\footnotesize$\abs{x_2-x_1}$};
  \draw[dashed] (B) -- (B |- A) node[midway,right]{\footnotesize$\abs{y_2-y_1}$};
  \draw (A) ++(0.5,0) arc (0:30:0.5);
  \node at ($(A)+(0.68,0.16)$) {\footnotesize$\theta$};
  \fill (A) circle (1pt) node[left]{\footnotesize$A(x_1,y_1)$};
  \fill (B) circle (1pt) node[above left]{\footnotesize$B(x_2,y_2)$};
\end{tikzpicture}
\end{center}

水平投影满足 $\abs{AB}\cos\theta=\abs{x_2-x_1}$，故 $\abs{AB}=\dfrac{1}{\abs{\cos\theta}}\abs{x_2-x_1}$。又由 $\sin^2\theta+\cos^2\theta=1$ 两边除以 $\cos^2\theta$：
\[
\tan^2\theta+1=\frac{1}{\cos^2\theta}
\quad\Longrightarrow\quad
\frac{1}{\abs{\cos\theta}}=\sqrt{k^2+1}.
\]

\state{弦长公式}{
    \[
    \abs{AB}=\sqrt{k^2+1}\;\abs{x_2-x_1}
    =\sqrt{k^2+1}\;\sqrt{(x_1+x_2)^2-4x_1x_2}.
    \]
}

把交点横坐标 $x_1,x_2$ 整体用韦达定理替换，就不用解出交点坐标了。

\ex{已知 $y_1=x^2+2x+3$ 与 $y_2=x+6$ 相交于 $A,B$ 两点，求线段 $AB$ 的长度。}

\sol{
    联立 $x^2+2x+3=x+6$，即 $x^2+x-3=0$。由韦达定理 $x_1+x_2=-1$、$x_1x_2=-3$，故
    \[
    (x_1-x_2)^2=(x_1+x_2)^2-4x_1x_2=1+12=13.
    \]
    直线 $y=x+6$ 斜率 $k=1$，代入弦长公式：
    \[
    \abs{AB}=\sqrt{k^2+1}\,\abs{x_2-x_1}=\sqrt2\cdot\sqrt{13}=\sqrt{26}.
    \]
}

\ex{已知圆 $x^2+y^2=4$ 与直线 $y=x+1$ 相交于 $A,B$ 两点，求线段 $AB$ 的长度。}

\sol{
    把 $y=x+1$ 代入 $x^2+y^2=4$：$x^2+(x+1)^2=4$，整理得 $2x^2+2x-3=0$。由韦达定理 $x_1+x_2=-1$、$x_1x_2=-\tfrac32$，故 $(x_1-x_2)^2=(x_1+x_2)^2-4x_1x_2=1+6=7$。斜率 $k=1$：
    \[
    \abs{AB}=\sqrt{k^2+1}\,\abs{x_2-x_1}=\sqrt2\cdot\sqrt7=\sqrt{14}.
    \]
}
